\section{A degree-316 ordinary-strong merger}
\label{sec:ordinary-merger}

\paragraph{Result and scope.}
There is an explicit equation in twelve strictly positive integer unknowns,
of exact total degree $316$, evaluated in $102$ binary arithmetic operations
($61$ multiplications and $41$ additions or subtractions), whose positive
unsolvability is equivalent in ZFC to $\operatorname{Con}(\mathrm{ZFC})$.
The three-equation system has residual degrees $2,26,288$ in $95$ operations,
or $94$ after the inherited positive transport-coordinate shift. A sum of
squares has degree $576$ in $103$ operations. Fixed numeral leaves and reuse
are free; multiplication by a numeral is charged.

This improves the authenticated degree-$897$, $108$-operation twelve-positive
construction on both measures. It does not reduce the eleven-positive
variable frontier, assert global optimality, prove consistency, or provide
a new end-to-end proof-assistant certificate. The native compiler and literal
K215/refutation calibration remain explicit inherited dependencies.

The structural change is to retain $f$ and eliminate the ordinary strong
coordinate $i$, together with $\tau$ and $y$. Thus the crucial polynomial is
$Q=\Delta(f^2-1)$ of degree $14$, not the normalized
$\Delta^2(ic^2)^2$ of degree $46$. A separated three-root sum recovers
rationality; rational Pell saturation recovers integrality. Removing a known
common denominator factor from its norm gives a degree-$288$ inner polynomial.

\subsection{Literal polynomial}
The supplied positive variables are
\[
 J,F,\alpha,z,f_0,T,s,w,\eta,\zeta,\gamma,t.
\]
Use exactly the entire authentic source31 Wang943 $121$-by-$8801$ coefficient
tuple $(\beta,K_0,M_C,M_F)$, with $\beta=2^{1064921}-1$ and $M_F$ already
the shifted mask. These are fixed integers, not unknowns. The payloads,
hashes, tuple adapter, and inherited ordinary proof are included in the
packet. Every other name below is a polynomial abbreviation:
\begin{align*}
q&=\beta J+1,& C_w&=q-F-\alpha,& X&=wq,&Y&=sq^3,\\
E&=XY,&k&=\eta+\zeta,&c&=kY+\eta,&a&=E+Y,\\
H&=4a+3,&\Delta&=a^2+H=(a+2)^2-1,&\ell&=EkY,\\
A&=1+\ell(\ell+k),&G&=1+\Delta c^2,&b_0&=X+ac.
\end{align*}
Define
\begin{align*}
R&=((q-1)(q-F)+\alpha)(q^2-1)+(M_C+qM_F)J,\\
N_t&=(K_0+w)C_w+q-F-z(q-1),\\
U&=4q^4,\qquad u=k-R-1,\\
D_*&=Ub_0+H(U\gamma-u),\qquad N_m=D_*^2-U^2G,\\
f&=f_0+1,\qquad Q=\Delta(f^2-1),\qquad d=Q-1,\\
V&=c(Tf-1)-Rf^2,\qquad E_a=QV^2-1,\qquad M_2=2f\ell,\\
Z&=c^4M_2^2 E_a,\qquad x=c^2t,\qquad r=c^4f^2A,\\
M&=d(x^2+Q+r)-Z,\\
C_0&=M^2+4d^2\{x^2(r-Q)-Qr\},\qquad B=2dQ-M,\\
N&=C_0^2-16d^2x^2rB^2.
\end{align*}
The equation is
\begin{equation}
 \boxed{\;N_t(1+N_m)(1+N)-1=0.\;}             \label{eq:ordinary316}
\end{equation}
The corresponding system is $N_t=1$, $N_m=0$, $N=0$. The sum of squares
is $(N_t-1)^2+N_m^2+N^2=0$.

The inner norm is exactly
\begin{equation}
 N=\prod_{\rho^2=Q,\;\tau^2=A}
 \bigl[d(c^2t-\rho-c^2f\tau)^2-c^4M_2^2E_a\bigr],
 \label{eq:ordinary-norm}
\end{equation}
with both formal signs of each root. To verify this identity, put
$b=c^2f\tau$, so $b^2=r$. Each factor is
$M-2dx\rho-2dxb+2d\rho b$. Its norm in $\rho$ is
$C_0+4dx(2dQ-M)b$; its norm in $b$ is the stated $N$.
This is a polynomial-ring identity, including degenerate specializations.
It does not assume independent Galois automorphisms after specialization.

For comparison, the eight-factor linear norm of
\[
 c^2dt-d\rho-c^2df\tau-c^2M_2\sqrt{dE_a}
\]
is $d^4N$. Eliminating the two signs of the last root proves this identity
and explains the denominator deflation. The actual circuit has no division.

\subsection{Unit guard and the main root}
Since $U$ is divisible by $4$,
\[
 N_m\equiv D_*^2\pmod {16},\qquad N\equiv C_0^2\pmod {16}.
\]
At a root of \eqref{eq:ordinary316}, all three factors are integer units.
Neither norm can equal $-2$, which is not a square modulo $4$. Hence
$N_m=N=0$ and $N_t=1$.

The authentic transport/mask interface then gives
\[
 q\ge16,\quad C_w>0,\quad 3q<R<q^4-2,\quad
 X\ge q,\quad Y\ge q^3,\quad E\ge q^4>R+2,
\]
\[
 k\ge2,\qquad kY<c<k(Y+1),\qquad a+2<2E.
\]
In particular $C_w\le0$ would imply
$N_t\le 1-F-(z-1)(q-1)\le0$, so it is impossible.
The shift $f=f_0+1$ ensures $f\ge2$, $Q>1$, and $d>0$;
no $f=1$ or zero-denominator exception is discarded.

From $N_m=0$, the rational root $g=D_*/U$ of $g^2=G$ is an integer.
It is positive: $H<5a$, $b_0>akY$, and $u<k$ give
\[
 D_*>Ub_0-Hk>ak(UY-5)>0.
\]
Reduction of $D_*=Ug$ modulo $U$ gives
\begin{equation}
 U\mid H(k-R-1).                                \label{eq:ordinary-absorber}
\end{equation}
Neither the omitted index quotient nor the old main quotient has yet
been recovered.

\subsection{Rationality from separated square classes}
Write $\rho=\sqrt Q>0$ and $\tau=\sqrt A>0$. If $V=0$, then $E_a=-1$
and every real factor in \eqref{eq:ordinary-norm} is strictly positive.
Thus $V\ne0$. Since $V$ is an integer, $E_a\ge Q-1=d$ and
\[
 y=\sqrt{E_a/d}\ge1.
\]
A vanishing factor gives signs $\epsilon_j\in\{1,-1\}$ with
\begin{equation}
 t=\epsilon_1\frac{\rho}{c^2}
    +\epsilon_2 f\tau+\epsilon_3 M_2y.           \label{eq:ordinary-sum}
\end{equation}
Before any rationality conclusion, one has
\[
 \rho<(a+2)f,\qquad \tau<\ell+k,\qquad
 \ell-k=k(EY-1)\ge2EY-2>2E>a+2.
\]
Consequently
\begin{equation}
 0<\frac\rho{c^2}<f\ell<f\tau,
 \qquad
 \frac\rho{c^2}+f\tau<2f\ell=M_2\le M_2y.
 \label{eq:ordinary-separation}
\end{equation}

Each term in \eqref{eq:ordinary-sum} is a rational multiple of a square
root of a positive integer; use $y=\sqrt{dE_a}/d$. Distinct squarefree
radicals are linearly independent over $\mathbb Q$. In a rational signed
sum, every irrational square class must cancel. The largest term exceeds
the sum of both others, so its class cannot cancel unless $y$ is rational.
Then the strict inequality between the two remaining terms excludes their
cancellation in an irrational class. Hence $\rho,\tau,y$ are rational.
In particular $\rho$ and $\tau$ are integers. This argument explicitly
allows coincident square classes.

\paragraph{Rational Pell saturation.}
For integers $m\ge2,V$ and rational $y$, the equation
\[
 m^2V^2-(m^2-1)y^2=1
\]
forces $y$ to be integral. Indeed, write $y=b/r$ in lowest terms. Then
$r^2\mid m^2-1$. For the positive nonsquare integer
$D_0=(m^2-1)/r^2$, the pairs $(m,r)$ and $(m|V|,|b|)$ are
integer-coordinate Pell solutions. If their fundamental-unit indices are
$j,l$, divisibility of their abscissas forces $l$ to be an odd multiple of
$j$: reduce modulo $2j$, reflect into $[0,j]$, and use strict growth with
$\chi_{n+2j}\equiv-\chi_n\pmod {\chi_j}$. Ordinate divisibility then gives
$r\mid b$, so $r=1$. The original equation excludes zero-coordinate cases.

Apply this lemma with $m=\rho$ to
\[
 \rho^2V^2-(\rho^2-1)y^2=1.
\]
Thus $y$ is integral. Equation \eqref{eq:ordinary-sum} now makes
$\rho/c^2$ integral. Put $i=\rho/c^2>0$. The recovered equations are
\[
 \tau^2=1+\ell(\ell+k),\qquad
 (ic^2)^2=\Delta(f^2-1),\qquad QV^2-(Q-1)y^2=1.
\]
All square-root signs may be made positive. No cancellation of $\Delta$
and no ordinary-rank assumption has occurred in this recovery.

\subsection{Ordered index recovery}
This part uses the preserved ordinary strong-rank theorem and the source50
absorber, with their hypotheses now verified. Let $A_0=a+2$ and
$P_0=2EY+1$. The two recovered Pell norms yield indices $p,n\ge1$ with
\[
 g=\chi_{A_0}(p),\quad c=\psi_{A_0}(p),\quad
 \tau=\chi_{P_0}(n),\quad k=2\psi_{P_0}(n).
\]
The first norm's fundamental integer-coordinate solution is
$(2EY+1,2)$ for discriminant $EY(EY+1)$. Since $P_0>A_0$ and $c>k$,
monotonicity gives $p>n$.

The Pell recurrence gives $k\equiv2n\pmod {4EY}$, and
$4EY=4ws^2q^7$ is divisible by $U$. Set $\mu=U/\gcd(U,H)$.
The number $H$ is odd; a prime common to $q,H$ must be $3$.
If $3\mid q$, then $H\equiv3\pmod {27}$, so its $3$-adic valuation
is exactly one. Hence
\[
 \gcd(U,H)\le3,\qquad \mu\ge\frac{4q^4}{3}>R+1.
\]
Together with \eqref{eq:ordinary-absorber} this gives
$2n=R+1+j\mu$, where positivity of $n$ forces $j\ge0$.
Thus $n\ge25$, $p\ge26$, and Pell growth proves
\[
 c>A_0\Delta^2,\qquad c>2p,\qquad c>2R.
\]

The ordinary strong-rank theorem now applies to
$(ic^2)^2=\Delta(f^2-1)$. It gives an index $m$ such that
\[
 f=\chi_{A_0}(m),\quad p\mid m,\quad c\mid m,\quad m>2p,
 \quad ic^2=\Delta\psi_{A_0}(m),\quad c^2\mid f^2-1.
\]
The theorem handles the possibly nonfundamental parameter in the underlying
squarefree Pell field. Its proof and the signed auxiliary argument are
included in the preserved ordinary source, Sections 4--5.

The odd auxiliary index $j_0$ satisfies $j_0\equiv\pm p\pmod m$ from
$V\equiv-c\pmod f$, and $j_0\equiv\pm R\pmod c$ from
$Q\equiv0\pmod c$, $V\equiv-R\pmod c$. Since $c\mid m$ and
$0<p,R<c/2$, it follows that $p=R$. The signed reduction uses
$2\chi_{A_0}(m-1)<\chi_{A_0}(m)$, not an assumed sign of $V$.

Now $n<R$, so $|2n-R-1|<R<\mu$ forces $2n=R+1$. Hence
$u=k-2n$ is divisible by $4EY$, and therefore by both $E,U$. Set
\[
 h=u/E,\qquad \gamma_{\rm old}=\gamma-u/U.
\]
The first is positive because $n\ge25$ and $\psi_{P_0}(n)>n$.
The main carrier gives $\gamma_{\rm old}=(g-b_0)/H$.
Since $n\ge2$, $k\ge4P_0>X$, so $c>X$ and
$g-b_0>(a+1)c-ac-X=c-X>0$.
All positive ordinary compiler coordinates are restored. Its inherited
soundness theorem extracts a blank K215 halt and a finite ZFC refutation.

\subsection{Completeness}
Conversely take a canonical positive ordinary compiler witness. Its index
identity is $p=R=2n-1$, so $u$ is a positive multiple of $4EY$ and of $U$.
Set
\[
 \gamma=\gamma_{\rm old}+u/U,\qquad f_0=f-1,
 \qquad t=i+f\tau+M_2y.
\]
These are positive integers. Then $D_*=Ug$, while the positive-sign
factor in \eqref{eq:ordinary-norm} is
$c^4M_2^2(dy^2-E_a)=0$. Thus $N_m=N=0$ and $N_t=1$.
The preserved change $i_{\rm ordinary}=\Delta i_{\rm normalized}$
provides ordinary witnesses from normalized ones when necessary.
This proves existence equivalence, hence the claimed consistency equivalence.

The shifted system uses $(K_0+w)C_w=F+z(q-1)$, corresponding to
$z_{\rm old}=z+1$. Canonical witnesses have $X=2^R$ and
$w>q>F$, which force $z_{\rm old}\ge2$. Thus the new $z$ is positive,
and the shift saves one operation.

\subsection{Exact degree and replay}
Put
\[
 q_0=\beta J,\quad Y_0=sq_0^3,\quad E_0=wq_0Y_0,
 \quad c_0=(\eta+\zeta)Y_0.
\]
These have degrees $1,4,6,5$. The inner leading forms are
\[
 Q_{14}=E_0^2f_0^2,\quad V_7=c_0Tf_0,\quad
 (M_2)_{12}=2f_0E_0c_0,\quad Z_{72}=4E_0^4c_0^8T^2f_0^6.
\]
All other terms of $M$ have degree at most $58$. Therefore
\[
 N_{288}=Z_{72}^4=2^8E_0^{16}c_0^{32}T^8f_0^{24}.
\]
The subtracted term in $N$ has degree at most $238$, so there is no
possible cancellation of this highest form.
For the main residual use
\[
 N_m=U^2\{(b_0+H\gamma)^2-G\}
      -2UHu(b_0+H\gamma)+H^2u^2.
\]
The bracket has degree $18$ and highest form $8E_0^2c_0\gamma$;
the other terms have degree at most $25,20$. Thus
\[
 (N_m)_{26}=128q_0^8E_0^2c_0\gamma.
\]
Since $(N_t)_2=w(q_0-F-\alpha)-zq_0$ is nonzero,
\[
 P_{316}=2^{15}q_0^8E_0^{18}c_0^{33}\gamma T^8f_0^{24}
             [w(q_0-F-\alpha)-zq_0]
\]
is nonzero for every authentic $\beta>0$. These are exact fixed-instance
degrees, not generic numerical tests.

The packet's independent arithmetic checker verifies $36$ exact signed
norm/deflation identities, $100$ modular identities, both residue guards,
$144$ pre-index bounds, and $24$ complete finite-field polynomial rays
including the actual fixed coefficients. It also verifies payload hashes,
analytic leaders, and transport shifts. Both ordinary Python and
\texttt{python -O} runs pass. The inert circuits have no dead gates or inputs.
Finite checks validate identities, counts and degrees, not the inherited
unbounded compiler theorem. The work is AI-assisted and makes no claim of
independent human review or kernel-certified consistency reduction.
